\documentclass[12pt]{article}
\usepackage[utf8]{inputenc}
\usepackage[T1]{fontenc}
\usepackage[margin=1in]{geometry}
\usepackage{setspace}
\usepackage{xcolor}
\usepackage[colorlinks=true, urlcolor=blue!60!black, linkcolor=blue!60!black]{hyperref}
\setlength{\parindent}{2em}
\setlength{\parskip}{1.25em}
\renewcommand{\baselinestretch}{1.0}
\title{\textit{Champion v. Ames}, 188 U.S. 321 (1903)}
\author{Case Brief}
\date{}

\begin{document}
\maketitle

\section{Case Metadata}
\begin{itemize}
  \item \textbf{Citation:} \textit{Champion v.\ Ames} (The Lottery Case), 188 U.S.\ 321 (1903) (23 S.\ Ct.\ 321; 47 L.\ Ed.\ 492)
  \item \textbf{Argued / Decided:} Argued Feb.\ 27--28, 1901; reargued Oct.\ 16--17, 1901; reargued before a full bench Dec.\ 15--16, 1902; decided February 23, 1903
  \item \textbf{Vote:} 5--4
  \item \textbf{Author:} Justice John Marshall Harlan (joined by Justices Henry Billings Brown, Edward Douglass White, Joseph McKenna, and Oliver Wendell Holmes Jr.)
  \item \textbf{Concurrences / Dissents:} No concurrences. Dissent by Chief Justice Melville W.\ Fuller (joined by Justices David J.\ Brewer, George Shiras Jr., and Rufus W.\ Peckham)
  \item \textbf{Chapter / Topic:} Federal Legislative Power --- Commerce Clause (Prohibition as Regulation; Federal ``Police Power'')
  \item \textbf{Primary Source URLs:} \href{https://supreme.justia.com/cases/federal/us/188/321/}{Justia U.S.\ Supreme Court Center} $|$ \href{https://www.law.cornell.edu/supremecourt/text/188/321}{Cornell LII}
\end{itemize}

\section{Facts and Procedural History}
In the late nineteenth century, Congress moved progressively to eradicate interstate lottery schemes after state-level prohibitions proved unable to stop out-of-state operators from using national transportation networks. Having already barred lottery materials from the federal mails and foreign importation, Congress enacted the Federal Lottery Act of March 2, 1895, ch.\ 191, \S~1, 28 Stat.\ 963, entitled ``An Act for the Suppression of Lottery Traffic through National and Interstate Commerce and the Postal Service, Subject to the Jurisdiction and Laws of the United States.'' Section~1 made it a federal crime to cause any lottery ticket, share, or advertisement to be ``carried from one state to another in the United States,'' punishable by up to two years' imprisonment, a \$1,000 fine, or both for a first offense.

On or about February 1, 1899, Charles F.\ Champion (alias W.\ W.\ Ogden), along with W.\ F.\ Champion and Charles B.\ Park, was indicted in the U.S.\ District Court for the Northern District of Texas for conspiracy under Revised Statutes \S~5440 to violate \S~1 of the 1895 Act. Specifically, the indictment alleged that the conspirators deposited a box containing two whole tickets of the Pan-American Lottery Company---which purported to conduct monthly drawings in Asunci\'{o}n, Paraguay, for a \$32,000 capital prize, selling whole tickets for \$2 and fractional tickets down to 25 cents---with the Wells-Fargo Express Company in Dallas, Texas, for shipment by rail to Fresno, California.

Champion was arrested in Chicago under a federal removal warrant to compel his appearance for trial in Texas. After declining to post bond, he was committed to the custody of U.S.\ Marshal John C.\ Ames and petitioned the U.S.\ Circuit Court for the Northern District of Illinois for a writ of habeas corpus, contending that \S~1 of the 1895 Act exceeded Congress's authority under the Commerce Clause. The Circuit Court dismissed the writ, and Champion appealed directly to the U.S.\ Supreme Court. The Court heard argument three separate times before resolving the case 5--4 after Justice Oliver Wendell Holmes Jr.\ joined the bench in December 1902.

\section{Issue(s) Presented}
\begin{enumerate}
  \item Does the interstate transportation of lottery tickets by an independent express carrier for hire constitute ``Commerce \dots\ among the several States'' within the meaning of Article~I, Section~8, Clause~3?
  \item Does Congress's power to ``regulate'' interstate commerce include the authority to \textit{prohibit} the interstate shipment of articles Congress deems injurious to public morals, or does such a prohibition unconstitutionally usurp the police power reserved to the states under the Tenth Amendment?
\end{enumerate}

\section{Holding and Rule}
Yes to both. Lottery tickets are subjects of traffic and therefore subjects of commerce, and their carriage from one state to another by independent carriers constitutes interstate commerce. \textbf{Rule:} Under the Commerce Clause, Congress's plenary authority to regulate interstate commerce includes the power to prohibit the interstate transportation of articles Congress deems injurious to public morals (here, lottery tickets), and doing so does not violate the Tenth Amendment or the Fifth Amendment's Due Process Clause. Judgment affirmed.

\section{Reasoning}
Writing for a five-Justice majority, Justice Harlan structured the Court's analysis in two sequential steps: first establishing that interstate carriage of lottery tickets is ``commerce,'' and second establishing that ``regulation'' encompasses ``prohibition.''

\begin{itemize}
  \item \textbf{Broad Definition of ``Commerce'':} Relying heavily on Chief Justice Marshall's foundational formulation in \textit{Gibbons v.\ Ogden}, 22 U.S.\ (9 Wheat.) 1 (1824), and subsequent nineteenth-century precedents (\textit{Brown v.\ Maryland}, \textit{The Passenger Cases}, \textit{Pensacola Telegraph Co.\ v.\ Western Union Telegraph Co.}, and \textit{Covington \& Cincinnati Bridge Co.\ v.\ Kentucky}), Harlan emphasized that ``commerce'' is not confined to the barter of tangible goods; it encompasses ``navigation, intercourse, communication, traffic, the transit of persons, and the transmission of messages by telegraph.'' Rejecting Champion's claim that lottery tickets lacked intrinsic value, Harlan noted that the tickets represented a right to draw a substantial monetary prize, commanded a market price among buyers and sellers, and were carried for hire by an interstate express company. Thus, they were ``subjects of traffic'' and therefore ``subjects of commerce.''
  \item \textbf{Regulation Includes Prohibition:} Invoking \textit{McCulloch v.\ Maryland}, 17 U.S.\ (4 Wheat.) 316 (1819), Harlan reasoned that Congress enjoys wide discretion over the means adapted to execute its enumerated powers. If Congress may regulate interstate lottery traffic, it is not forced merely to prescribe the manner of shipment while tolerating an enterprise long condemned by the Court itself as a ``widespread pestilence'' (\textit{Phalen v.\ Virginia}, 49 U.S.\ (8 How.) 163, 168 (1850)). Harlan analogized the lottery ban to three accepted prohibitory uses of the Commerce Clause: (1) Congress's ban on transporting diseased cattle across state lines (\textit{Reid v.\ Colorado}, 187 U.S.\ 137 (1902)); (2) the Sherman Antitrust Act's prohibition of contracts restraining interstate trade (\textit{Addyston Pipe \& Steel Co.\ v.\ United States}, 175 U.S.\ 211 (1899)); and (3) the Wilson Act's effective exclusion of original-package liquor shipments into dry states (\textit{In re Rahrer}, 140 U.S.\ 545 (1891)).
  \item \textbf{No Violation of the Tenth or Fifth Amendments:} Harlan dismissed the Tenth Amendment objection directly: because the power to regulate interstate commerce is expressly delegated to Congress, exercising it cannot invade powers reserved to the states. Moreover, Congress had not regulated purely intrastate lottery sales; it had merely supplemented state anti-lottery policies by preventing interstate channels from defeating them. As for Fifth Amendment liberty of contract (\textit{Allgeyer v.\ Louisiana}, 165 U.S.\ 578 (1897)), no person has a constitutional right to inject into interstate commerce an article ``confessedly injurious to the public morals.''
  \item \textbf{Political Safeguards Against Abuse:} Confronting the slippery-slope objection that Congress might arbitrarily ban harmless or valuable goods from interstate commerce, Harlan declined to map the outer boundaries of the Commerce Clause in advance. Quoting \textit{Gibbons}, he underscored that the possible abuse of a power does not disprove its existence, and that for unwise but authorized legislation, the primary check lies in elections and the political process.
\end{itemize}

\section{Concurring and Dissenting Opinions}
\textbf{Chief Justice Fuller (dissenting, joined by Justices Brewer, Shiras, and Peckham):} Fuller delivered a vigorous dual-federalism dissent attacking both prongs of Harlan's opinion.
\begin{enumerate}
  \item \textbf{Pretext and the Tenth Amendment:} Fuller argued that the Act's avowed purpose---stated in its very title---was to suppress lotteries, a quintessentially local moral regulation belonging exclusively to the states' police power (\textit{United States v.\ E.\ C.\ Knight Co.}, 156 U.S.\ 1 (1895)). While Congress could exclude lottery matter from the federal mails under the Postal Clause (\textit{Ex parte Jackson}, 96 U.S.\ 727 (1878); \textit{In re Rapier}, 143 U.S.\ 110 (1892)), \textit{Jackson} had expressly disclaimed federal authority to ban non-postal carriage. Invoking \textit{McCulloch}'s ``pretext'' warning, Fuller warned that upholding the statute endowed Congress with a general national police power.
  \item \textbf{Lottery Tickets Are Not Articles of Commerce:} Analogizing lottery tickets to insurance contracts---which the Court had held were personal contracts rather than articles of commerce in \textit{Paul v.\ Virginia}, 75 U.S.\ (8 Wall.) 168 (1869)---Fuller contended that a noncommercial piece of paper does not magically become an ``article of commerce'' merely because it is placed in an express package. If mere carriage sufficed, ``an invitation to dine, or to take a drive, or a note of introduction'' would become interstate commerce.
  \item \textbf{Asymmetry Between Foreign and Interstate Commerce:} Fuller argued that Congress's power over interstate commerce was granted ``to form a more perfect union by freeing such commerce from state discrimination, and not to transfer the power of restriction.'' Unlike foreign commerce (where a sovereign nation may impose embargoes) or Indian commerce, interstate commerce is constrained by Article~IV's Privileges and Immunities Clause and the structural preservation of state autonomy. Diseased cattle could be excluded only because they were physically infectious to other articles in transit; lottery tickets ``cannot communicate bad principles by contact.''
\end{enumerate}

\section{Key Quotations}
\begin{enumerate}
  \item ``We are of opinion that lottery tickets are subjects of traffic, and therefore are subjects of commerce, and the regulation of the carriage of such tickets from state to state, at least by independent carriers, is a regulation of commerce among the several states.'' --- Justice Harlan, 188 U.S.\ at 354 (\href{https://supreme.justia.com/cases/federal/us/188/321/}{Justia}).
  \item ``If a state, when considering legislation for the suppression of lotteries within its own limits, may properly take into view the evils that inhere in the raising of money in that mode, why may not Congress, invested with the power to regulate commerce among the several states, provide that such commerce shall not be polluted by the carrying of lottery tickets from one state to another?'' --- Justice Harlan, 188 U.S.\ at 356 (\href{https://supreme.justia.com/cases/federal/us/188/321/}{Justia}).
  \item ``As a state may, for the purpose of guarding the morals of its own people, forbid all sales of lottery tickets within its limits, so Congress, for the purpose of guarding the people of the United States against the `widespread pestilence of lotteries' and to protect the commerce which concerns all the states, may prohibit the carrying of lottery tickets from one state to another.'' --- Justice Harlan, 188 U.S.\ at 357 (\href{https://supreme.justia.com/cases/federal/us/188/321/}{Justia}).
  \item ``But, as often said, the possible abuse of a power is not an argument against its existence. There is probably no governmental power that may not be exerted to the injury of the public.'' --- Justice Harlan, 188 U.S.\ at 363 (\href{https://supreme.justia.com/cases/federal/us/188/321/}{Justia}).
  \item ``An invitation to dine, or to take a drive, or a note of introduction, all become articles of commerce under the ruling in this case by being deposited with an express company for transportation. \dots\ It is a long step in the direction of wiping out all traces of state lines, and the creation of a centralized government.'' --- Chief Justice Fuller (dissenting), 188 U.S.\ at 371 (\href{https://supreme.justia.com/cases/federal/us/188/321/}{Justia}).
\end{enumerate}

\section{Doctrinal Significance and Modern Legacy}
\textit{Champion v.\ Ames} is the foundational precedent establishing a de facto federal ``police power'' exercised through the channels of interstate commerce. Immediately following \textit{Champion}, the Court relied on its commerce-prohibiting principle to uphold the Pure Food and Drug Act in \href{https://supreme.justia.com/cases/federal/us/220/45/}{\textit{Hipolite Egg Co.\ v.\ United States}}, 220 U.S.\ 45 (1911), the Mann Act (White-Slave Traffic Act) in \href{https://supreme.justia.com/cases/federal/us/227/308/}{\textit{Hoke v.\ United States}}, 227 U.S.\ 308 (1913), and the National Motor Vehicle Theft Act in \href{https://supreme.justia.com/cases/federal/us/267/432/}{\textit{Brooks v.\ United States}}, 267 U.S.\ 432 (1925).

The Court briefly attempted to cabin \textit{Champion} during the \textit{Lochner} era in \href{https://supreme.justia.com/cases/federal/us/247/251/}{\textit{Hammer v.\ Dagenhart}}, 247 U.S.\ 251 (1918), striking down the federal Child Labor Act by distinguishing goods ``harmful in themselves'' (like lottery tickets or impure food) from ``harmless'' manufactured goods produced by child labor. That distinction proved unworkable---after eroding it in \href{https://supreme.justia.com/cases/federal/us/299/334/}{\textit{Kentucky Whip \& Collar Co.\ v.\ Illinois Central R.\ Co.}}, 299 U.S.\ 334 (1937) (prison-made goods), a unanimous Court overruled \textit{Hammer} in \href{https://supreme.justia.com/cases/federal/us/312/100/}{\textit{United States v.\ Darby}}, 312 U.S.\ 100 (1941), fully vindicating \textit{Champion} and Justice Holmes's \textit{Hammer} dissent by holding that Congress's motive in barring goods from interstate shipment is constitutionally irrelevant. \textit{Champion}'s moral-purpose principle likewise underpinned Title~II of the Civil Rights Act of 1964 in \href{https://supreme.justia.com/cases/federal/us/379/241/}{\textit{Heart of Atlanta Motel, Inc.\ v.\ United States}}, 379 U.S.\ 241 (1964).

Through 2026, \textit{Champion} remains undisturbed bedrock law as the first pillar of the modern three-category Commerce Clause framework articulated in \href{https://supreme.justia.com/cases/federal/us/514/549/}{\textit{United States v.\ Lopez}}, 514 U.S.\ 549, 558 (1995): Congress's plenary authority to regulate ``the use of the channels of interstate commerce.'' Even as the modern Court has reasserted outer limits on federal legislative power---requiring economic activity under the ``substantial effects'' prong in \textit{Lopez} and \href{https://supreme.justia.com/cases/federal/us/529/598/case.pdf}{\textit{United States v.\ Morrison}}, 529 U.S.\ 598 (2000), prohibiting Congress from compelling commercial inactivity in \href{https://supreme.justia.com/cases/federal/us/567/519/}{\textit{NFIB v.\ Sebelius}}, 567 U.S.\ 519 (2012), and deploying federalism clear-statement canons to prevent treaty-implementing or federal criminal statutes from swallowing local assaults in \href{https://supreme.justia.com/cases/federal/us/572/844/}{\textit{Bond v.\ United States}}, 572 U.S.\ 844 (2014)---the Court has never retreated from \textit{Champion}'s core rule that Congress may close the channels of interstate commerce to disfavored articles, whether firearms, controlled substances (\href{https://supreme.justia.com/cases/federal/us/545/1/}{\textit{Gonzales v.\ Raich}}, 545 U.S.\ 1 (2005)), or digital contraband.

\section{Pedagogical Critiques and Exam Hypotheticals}
\textbf{Pedagogical Critiques:}
\begin{itemize}
  \item \textbf{The ``Pollution'' Metaphor and Scope Creep:} Harlan leaned heavily on the morally noxious character of lotteries, suggesting commerce could be protected from ``pollution.'' Yet as Fuller foresaw, once Congress may prohibit interstate movement based on its own moral or policy disapproval, the ``harmful-in-itself'' limitation collapses---producing the doctrinal whiplash between \textit{Champion}, \textit{Hammer}, and \textit{Darby}.
  \item \textbf{Formalism vs.\ Pretext:} Harlan's insistence that the 1895 Act left ``completely internal'' state affairs untouched is formally accurate but functionally illusory. Because national distribution is the lifeblood of most commercial enterprises, banning interstate shipment effectively lets Congress regulate upstream production and downstream local behavior without ever needing to claim a formal police power.
\end{itemize}

\textbf{Exam Hypothetical:} Congress enacts the ``Interstate Reproductive Information and Materials Act,'' making it a federal felony to ship any over-the-counter abortifacient medication or printed instructional brochure across state lines via private courier. A distributor in State~A (where the medication is legal) ships a package via FedEx to a clinic in State~B (where it is also legal) and challenges the indictment under the Commerce Clause and Tenth Amendment.
\begin{itemize}
  \item \textbf{Government's Best Argument (Channel-Closing Under \textit{Champion} and \textit{Darby}):} Under \textit{Champion} and \textit{Darby}, Congress's authority over the channels of interstate commerce is plenary and includes the power to prohibit interstate shipment of any article, regardless of whether the underlying motive is moral or regulatory, and even if both the origin and destination states permit the article locally.
  \item \textbf{Challenger's Best Argument (Opposite-Side Contrast):} First, distinguish \textit{Champion} on external constitutional limits: Harlan expressly acknowledged that the commerce power ``may not be exercised so as to infringe rights secured or protected by that instrument'' (188 U.S.\ at 362--363); banning the interstate shipment of \textit{printed brochures} directly implicates the First Amendment (an issue Fuller flagged as absent in \textit{Champion}). Second, unlike \textit{Champion}, where Harlan stressed that Congress was ``supplement[ing] the action of those states---perhaps all of them'' that banned lotteries, here Congress is overriding the permissive policies of both State~A and State~B, reviving Fuller's structural objection that Congress is using channel-closing as a pretextual national police power. (Under current doctrine, the First Amendment challenge to the brochures is far stronger than the Tenth Amendment/Commerce Clause challenge to the physical pills.)
\end{itemize}

\end{document}
